\documentclass[11pt,a4paper]{article}
\usepackage[T2A,T1]{fontenc}
\usepackage[utf8]{inputenc}
\usepackage[english,russian]{babel}
\usepackage[a4paper,margin=22mm]{geometry}
\usepackage{microtype}
\usepackage{lmodern}
\usepackage{parskip}
\usepackage{enumitem}
\usepackage{booktabs,longtable,tabularx,array}
\usepackage{xcolor}
\usepackage{hyperref}
\usepackage{fancyhdr}
\usepackage{titlesec}
\usepackage{listings}
\usepackage{tcolorbox}
\usepackage{lastpage}
\usepackage{pifont}
\usepackage{amsmath,amssymb}
\usepackage{graphicx}
\usepackage{needspace}
\usepackage{etoolbox}

\definecolor{navy}{HTML}{17324D}
\definecolor{blue}{HTML}{2F6690}
\definecolor{lightblue}{HTML}{EDF4F8}
\definecolor{green}{HTML}{2E6F40}
\definecolor{lightgreen}{HTML}{EDF7F0}
\definecolor{red}{HTML}{9E2A2B}
\definecolor{lightred}{HTML}{FBEFEF}
\definecolor{amber}{HTML}{8A5A00}
\definecolor{lightamber}{HTML}{FFF7E6}
\definecolor{gray}{HTML}{5F6B73}
\definecolor{lightgray}{HTML}{F4F6F7}

\hypersetup{colorlinks=true,linkcolor=navy,urlcolor=blue,citecolor=blue,pdftitle={Math CI/CD: Human Operating Manual},pdfauthor={Math CI/CD Project}}
\pagestyle{fancy}
\fancyhf{}
\lhead{\textcolor{gray}{Math CI/CD Human Manual}}
\rhead{\textcolor{gray}{v0.4.3}}
\cfoot{\textcolor{gray}{\thepage/\pageref{LastPage}}}
\renewcommand{\headrulewidth}{0.3pt}
\setlength{\headheight}{14pt}
\titleformat{\section}{\Large\bfseries\color{navy}}{\thesection}{0.7em}{}
\titleformat{\subsection}{\large\bfseries\color{blue}}{\thesubsection}{0.7em}{}
\titleformat{\subsubsection}{\normalsize\bfseries\color{navy}}{\thesubsubsection}{0.7em}{}
\setlist[itemize]{leftmargin=1.5em,itemsep=2pt,topsep=3pt}
\setlist[enumerate]{leftmargin=1.7em,itemsep=3pt,topsep=3pt}
\renewcommand{\arraystretch}{1.18}
\newcolumntype{Y}{>{\raggedright\arraybackslash}X}
\newcommand{\status}[1]{\texttt{#1}}
\newcommand{\cmd}[1]{\texttt{#1}}
\newcommand{\ok}{\textcolor{green}{\ding{51}}}
\newcommand{\bad}{\textcolor{red}{\ding{55}}}

\lstdefinestyle{shell}{basicstyle=\ttfamily\small,backgroundcolor=\color{lightgray},frame=single,rulecolor=\color{gray!35},breaklines=true,columns=fullflexible,keepspaces=true,showstringspaces=false,xleftmargin=4pt,xrightmargin=4pt,aboveskip=7pt,belowskip=7pt}
\lstdefinestyle{lean}{style=shell}
\lstdefinestyle{yaml}{style=shell}
\lstset{style=shell}

\newtcolorbox{principle}[1][]{colback=lightblue,colframe=blue,title=\textbf{Principle},fonttitle=\bfseries,boxrule=0.6pt,arc=2pt,#1}
\newtcolorbox{warningbox}[1][]{colback=lightamber,colframe=amber,title=\textbf{Attention},fonttitle=\bfseries,boxrule=0.6pt,arc=2pt,#1}
\newtcolorbox{dangerbox}[1][]{colback=lightred,colframe=red,title=\textbf{Forbidden / blocking condition},fonttitle=\bfseries,boxrule=0.6pt,arc=2pt,#1}
\newtcolorbox{checklistbox}[1][]{colback=lightgreen,colframe=green,title=\textbf{Checklist},fonttitle=\bfseries,boxrule=0.6pt,arc=2pt,#1}

\begin{document}
\selectlanguage{russian}
\begin{titlepage}
\centering
\vspace*{24mm}
{\Huge\bfseries\color{navy} Math CI/CD\par}
\vspace{5mm}
{\LARGE Руководство по AI-ассистированным математическим исследованиям\par}
\vspace{9mm}
{\large Канонический human operating manual / SOP\par}
\vfill
\begin{tcolorbox}[width=.82\textwidth,colback=lightblue,colframe=blue,arc=2pt]
\centering
\textbf{Назначение:} сделать математическую разработку с Lean и LLM воспроизводимой, проверяемой, понятной человеку и устойчивой к незаметному изменению утверждения, скрытым аксиомам, устаревшим review и ложным заявлениям о новизне.
\end{tcolorbox}
\vfill
{\large Версия 0.4.3\par}
{\normalsize 27 сентября 2026\par}
\vspace{8mm}
{\small Нормативный источник для людей: \texttt{docs/HUMAN\_MANUAL.tex}. Markdown-файлы репозитория являются краткими operational notes и не заменяют этот документ.\par}
\end{titlepage}

\tableofcontents
\clearpage

\section{Назначение и модель доверия}

Этот репозиторий является не только кодовой базой, но и \textbf{аудиторским следом математического результата}. Работа считается завершённой не тогда, когда агент написал Lean-код и не тогда, когда \cmd{lake build} зелёный, а когда раздельно собраны и актуальны доказательства пяти разных утверждений.

\begin{principle}
Никогда не объединяйте в один вопрос:
\begin{enumerate}
\item принимает ли Lean proof term;
\item соответствует ли Lean statement математике, которую команда намеревалась доказать;
\item понимает ли человек ключевые переходы доказательства;
\item является ли результат новым относительно литературы;
\item воспроизводится ли результат независимо из чистого checkout.
\end{enumerate}
Зелёный Lean build отвечает только на первый вопрос относительно закодированных определений и доверенной базы.
\end{principle}

Финальный статус \status{VERIFIED\_RESULT} не присваивается LLM и не эквивалентен успешной компиляции. Это состояние governance, которое допускается только после согласования machine evidence и требуемых human reviews.

\subsection{Что автоматизация может и не может гарантировать}

\begin{tabularx}{\textwidth}{@{}p{.29\textwidth}YY@{}}
\toprule
Слой & Автоматизация хорошо проверяет & Требует содержательной человеческой проверки \\
\midrule
Lean/kernel & build, placeholders, зарегистрированные declarations, зависимости, trusted-base policy & адекватность определения и научный смысл формализации \\
Statement control & frozen hash, semantic diff, drift & является ли новая формулировка правильной постановкой задачи \\
Proof evidence & fingerprints, stale review, dependency surface & понимание нетривиальных переходов и математическая ценность \\
Literature & наличие воспроизводимого журнала, provenance fields & полнота поиска, эквивалентность/приоритет, корректность novelty language \\
Reproduction & commit/environment/evidence binding & независимость reproducer и отсутствие скрытого generator context \\
\bottomrule
\end{tabularx}

\section{Роли и разделение ответственности}

Один человек может выполнять несколько ролей в маленькой группе, но роли должны оставаться концептуально раздельными. Нельзя превращать авторство proof в автоматическое semantic, novelty или final approval.

\begin{longtable}{@{}p{.22\textwidth}p{.72\textwidth}@{}}
\toprule
\textbf{Роль} & \textbf{Ответственность} \\
\midrule
\endhead
Research Owner & Формулирует научный вопрос, область применимости, intended definitions и interpretation. Регистрирует гипотезу до длительного proof search. Не усиливает assumptions молча, когда proof search затрудняется. \\
Formalizer & Переводит intended statement в Lean, поддерживает buildability, документирует нетривиальные formalization choices. Может предложить изменение statement, но не менять frozen mathematics скрытно. \\
Skeptic / adversarial reviewer & Ищет контрпримеры, вырожденные случаи, лишние assumptions, circularity, tautological definitions, более сильные известные результаты и пути тривиализации. \\
Proof-understanding reviewer & Восстанавливает обычное математическое доказательство, выделяет ключевые переходы и проверяет, что proof не только существует как term, но понятен. \\
Literature / provenance reviewer & Пытается опровергнуть novelty; фиксирует запросы, индексы, кандидатов, citation chasing, equivalence checks и proof fingerprints. \\
Independent reproducer & Работает из чистого checkout на фиксированном commit/environment, желательно без приватного контекста генератора. \\
Final reviewer & Проверяет соответствие claim language накопленному evidence и отсутствие stale reviews. \\
LLM / agents & Могут предлагать conjectures, Lean code, counterexamples, queries, dependency summaries и critiques. Не имеют права имитировать human sign-off, объявлять novelty доказанной или менять frozen statement без явного SCR. \\
\bottomrule
\end{longtable}

\section{Исследовательский объект и структура репозитория}

Каждое существенное утверждение получает неизменяемый идентификатор вида \status{H-0042}. Один H-ID соответствует одному научному объекту. Материально другое утверждение получает новый H-ID, а связь фиксируется в metadata.

\begin{lstlisting}
researchctl new H-0042 "Scaling rigidity candidate"
git switch -c hyp/H-0042-scaling-rigidity
\end{lstlisting}

Каноническая карточка: \texttt{hypotheses/H-0042.yaml}. Связанное evidence хранится под тем же ID:

\begin{lstlisting}
lit/H-0042.md
proofs/H-0042.md
provenance/H-0042.yaml
ablations/H-0042/
reviews/H-0042/
reports/.../H-0042.*
\end{lstlisting}

\begin{warningbox}
Markdown в \texttt{lit/} и \texttt{proofs/} допустим как рабочий журнал и краткая реконструкция. Нормативные инструкции для людей должны жить в LaTeX-документации. Machine-readable состояние хранится в YAML/JSON; Lean остаётся источником формальной истины.
\end{warningbox}

\section{Жизненный цикл результата}

Нормальный путь:

\begin{center}
\status{CANDIDATE} $\rightarrow$ \status{SPEC\_DRAFT} $\rightarrow$ \status{SPEC\_FROZEN} $\rightarrow$ \status{FORMALIZED} $\rightarrow$ \status{SEMANTIC\_REVIEWED} $\rightarrow$ \status{PROOF\_RECONSTRUCTED} $\rightarrow$ \status{NOVELTY\_REVIEWED} $\rightarrow$ \status{INDEPENDENTLY\_REPRODUCED} $\rightarrow$ \status{VERIFIED\_RESULT}
\end{center}

\begin{longtable}{@{}p{.25\textwidth}p{.67\textwidth}@{}}
\toprule
\textbf{Статус} & \textbf{Научный смысл} \\
\midrule
\endhead
\status{CANDIDATE} & Есть вопрос/гипотеза; точная формализация может быть неполной. Фиксируются intended phenomenon, scope, assumptions, известные связи и первые falsification targets. \\
\status{SPEC\_DRAFT} & Стабилизируются definitions, quantifiers, domains и assumptions. Это правильное место для semantic redesign. \\
\status{SPEC\_FROZEN} & Точный Lean statement заморожен и получает SHA-256 snapshot. Изменение binders/assumptions/conclusion после этого является scientific change, а не refactor. \\
\status{FORMALIZED} & Полный Lean proof существует без запрещённых placeholders и с допустимой trusted base. Не означает meaningful/novel/understood. \\
\status{SEMANTIC\_REVIEWED} & Человек проверил соответствие Lean statement intended mathematics и отсутствие скрытой тривиализации. \\
\status{PROOF\_RECONSTRUCTED} & Человек понимает ключевые переходы, проверил circularity/trivialization risks и review surface. \\
\status{NOVELTY\_REVIEWED} & Literature/provenance audit достаточно полон для заявляемого уровня novelty; поиск воспроизводим. \\
\status{INDEPENDENTLY\_REPRODUCED} & Отдельный reproducer воспроизвёл результат из чистого checkout/environment и связал review с конкретными fingerprints. \\
\status{VERIFIED\_RESULT} & Все требуемые evidence свежи, sign-offs выполнены, policy проходит, final claim language утверждён. \\
\bottomrule
\end{longtable}

Допустимые terminal/side outcomes: \status{REFUTED}, \status{KNOWN\_RESULT}, \status{SPECIFICATION\_BUG}, \status{BLOCKED}, \status{INCONCLUSIVE}. Полезные отрицательные результаты следует сохранять и merge'ить, а не удалять.

\section{Git и GitHub: рабочий протокол}

\subsection{Ветки}

\begin{lstlisting}
hyp/H-0042-short-name       # hypothesis/proof development
formal/H-0042-short-name    # formalization
skeptic/H-0042-short-name   # counterexamples/adversarial work
lit/H-0042-short-name       # literature/provenance
proof/H-0042-short-name     # human proof reconstruction
scr/H-0042-short-name       # statement change request
infra/short-name            # CI/CD infrastructure
\end{lstlisting}

Не используйте долгоживущие персональные ветки. H-ID должен присутствовать в branch/PR, изменяющем scientific object.

\subsection{Issue и Draft PR}

Issue создаётся до устойчивой работы и содержит H-ID, informal claim, scope, assumptions/definitions, мотивацию, known related results, counterexample targets, owners/reviewers и текущий статус. Draft PR открывается рано и становится общей coordination surface для людей и агентов.

Каждый серьёзный PR должен отвечать: что изменилось математически; frozen ли statement; какие evidence изменились; какие reviews стали stale; что доказано/опровергнуто; что искалось в литературе; какие checks необходимо повторить.

\subsection{Коммиты}

Предпочтительны маленькие аудируемые commits:

\begin{lstlisting}
H-0042: register scaling rigidity candidate
H-0042: freeze theorem statement
H-0042: prove compactness bridge lemma
H-0042: add counterexample search results
H-0042: record nearest prior-art candidates
H-0042: reconstruct proof skeleton
infra: reject invalid Lean references at candidate stage
\end{lstlisting}

Не смешивайте изменение CI/CD и headline mathematical proof в одном commit без необходимости.

\section{Statement freeze и управление изменениями}

Перед freeze три Lean reference fields должны быть различены:

\begin{lstlisting}[style=yaml]
statement:
  lean_file: MathCI/Scaling.lean
  lean_module: MathCI.Scaling
  lean_declaration: MathCI.scaling_rigidity
\end{lstlisting}

\texttt{lean\_file} --- путь файловой системы; \texttt{lean\_module} --- импортируемый Lean module; \texttt{lean\_declaration} --- fully-qualified theorem/definition.

\begin{lstlisting}
researchctl validate
researchctl freeze H-0042
researchctl semantic-diff H-0042
\end{lstlisting}

После freeze proof refactoring разрешён при неизменном statement. Любое semantic изменение требует явного Statement Change Request: вернуть объект в \status{SPEC\_DRAFT}, объяснить причину, заморозить новую версию и повторить downstream reviews.

\begin{dangerbox}
Запрещено менять assumption, domain, binder, conclusion или definition только для того, чтобы агент смог завершить proof, а затем выдавать это за proof исходной гипотезы.
\end{dangerbox}

\section{Формализация, trusted base и зависимости}

Во время Lean-разработки запускайте:

\begin{lstlisting}
lake build
researchctl validate
researchctl semantic-diff H-0042 --fail-on-change
\end{lstlisting}

Перед \status{FORMALIZED}:

\begin{lstlisting}
researchctl audit H-0042
researchctl axioms-all
researchctl deps-exact H-0042
researchctl review-surface H-0042
\end{lstlisting}

\subsection{Placeholder proofs}

Lean может принять \texttt{sorry} с warning. Math CI/CD обязан трактовать \texttt{sorry}/\texttt{admit} как hard failure. Успешный \cmd{lake build} сам по себе не означает complete proof.

\subsection{Axiom policy}

Trusted base задаётся allowlist policy. Project-specific axiom является научным предположением и должен быть явно обоснован. Нельзя добавлять axiom в allowlist только для того, чтобы CI стал зелёным.

Fresh transitive traversal используется как основной gate; \texttt{\#print axioms} сохраняется как диагностическое сравнение.

\subsection{Exact dependency graph и review surface}

Exact graph строится из elaborated Lean environment. Он помогает найти project lemmas, на которых концентрируется человеческий review. Граф не заменяет математическое понимание и не является оценкой научной важности.

\section{Semantic review}

После готовности evidence:

\begin{lstlisting}
researchctl fingerprint H-0042
researchctl review-init H-0042
\end{lstlisting}

Semantic reviewer обязан проверить definitions, assumptions, quantifiers/domains/boundaries, conclusion, vacuous/degenerate cases и отсутствие скрытого усиления спецификации. Review связан с frozen statement hash; изменение statement делает его stale/invalid.

\begin{checklistbox}
Перед \status{SEMANTIC\_REVIEWED}: definitions соответствуют intended objects; assumptions явно перечислены; conclusion совпадает с human claim; degenerate cases рассмотрены; theorem не стал тривиальным из-за encoding; reviewer не является LLM-подписью.
\end{checklistbox}

\section{Понимание proof и реконструкция}

\texttt{proofs/H-0042.md} является рабочей human-readable реконструкцией, а не line-by-line переводом Lean. Она должна содержать conventional statement, 2--10 ключевых lemmas/transitions, границу routine library machinery и project-specific reasoning, использование assumptions и потенциальные circularity/trivialization points.

\begin{lstlisting}
researchctl deps-exact H-0042
researchctl provenance-init H-0042
researchctl review-surface H-0042
\end{lstlisting}

Proof review привязывается к formalization fingerprint. Изменение proof source, proof reconstruction или relevant environment делает старый review stale.

\section{Assumption ablation и falsification}

\begin{lstlisting}
researchctl ablate H-0042
\end{lstlisting}

Для каждого содержательного assumption попытайтесь удалить/ослабить его, построить counterexample без него или указать точный proof step, где он необходим. Нельзя называть assumption необходимым только потому, что текущий proof его использует.

Adversarial review должен проверять empty domains, zero/one limits, symmetry-degenerate и boundary cases, equivalent reparameterizations, tautological definitions и stronger known results.

\section{Литературный обзор и provenance}

\begin{principle}
Цель literature review --- не подтвердить novelty, а активно попытаться её опровергнуть.
\end{principle}

\subsection{Обязательные семейства поиска}

Ищите exact terminology; стандартные термины вместо project names; equivalent formulations; stronger results; weaker/general/special cases; characteristic proof ingredients; необычную последовательность reductions (proof fingerprint); источники major imported/library lemmas.

Для каждого запроса фиксируйте дату, database/index/source, воспроизводимый query, candidate results и классификацию каждого серьёзного кандидата как equivalent/stronger/weaker/orthogonal/unresolved. Для сильных кандидатов делайте backward и forward citation chasing, когда применимо.

\subsection{Provenance categories}

\status{PROJECT\_NEW}, \status{PROJECT\_PRIOR}, \status{MATHLIB}, \status{EXTERNAL\_PUBLISHED}, \status{STANDARD\_FACT}, \status{UNKNOWN}. \status{UNKNOWN} допустим во время исследования, но не должен незаметно пережить сильный novelty claim.

Agents могут сообщать: ``matching result не найден в перечисленных источниках''. Они не должны превращать это в ``theorem new''. При найденном эквивалентном prior result используйте \status{KNOWN\_RESULT}.

\section{Freshness reviews и fingerprints}

\begin{lstlisting}
researchctl report H-0042
researchctl fingerprint H-0042
\end{lstlisting}

Report должен различать \texttt{PASS / VALID}, \texttt{PASS / STALE}, \texttt{MISSING}, \texttt{INCONCLUSIVE}. \texttt{STALE} означает, что approval относился к старой версии evidence.

\begin{dangerbox}
Никогда не ``освежайте'' \texttt{\detokenize{*_sha256_reviewed}} механически. Правильное действие при STALE --- повторно проверить изменившийся материал и только затем записать новый review hash.
\end{dangerbox}

Fingerprint domains должны различать statement, formalization, literature/provenance, environment и complete evidence. Это позволяет повторять только затронутые проверки, не стирая независимые approvals без причины.

\section{Independent reproduction}

Reproducer предпочтительно работает в чистом clone и без generator chat history:

\begin{lstlisting}
git clone <repo> clean-reproduction
cd clean-reproduction
git checkout <commit>
python -m pip install -e '.[dev]'
lake build
researchctl validate
researchctl axioms-all
researchctl deps-exact H-0042
researchctl env-lock
researchctl audit H-0042
\end{lstlisting}

Фиксируются exact commit, environment и способность восстановить human proof без скрытого контекста генератора.

\section{CI: обязательные проверки и их смысл}

\begin{longtable}{@{}p{.22\textwidth}p{.70\textwidth}@{}}
\toprule
\textbf{Check} & \textbf{Что он означает} \\
\midrule
\endhead
\texttt{protocol-gate} & Metadata, valid Lean references, statement freeze, status requirements, freshness требуемых reviews, placeholders, provenance и final policy. \\
\texttt{unit-tests} & Тестирует реализацию \texttt{researchctl}. Не тестирует вашу математику напрямую. \\
\texttt{acceptance-suite} & Запускает sabotage regressions: baseline, \texttt{sorry}, invalid Lean references, statement drift, stale proof review и другие control failures. \\
\texttt{agent-protocol} & Проверяет полноту canonical prompt files, вычисляет общий/role SHA-256 и публикует manifest как CI artifact. \\
\texttt{blind-context} & Строит изолированные role-specific context packets, проверяет visibility profiles и сканирует пакеты на утечки запрещённого глобального контекста. \\
\texttt{workflow-audit} & Проверяет опасные GitHub Actions patterns. Privileged execution untrusted PR code с secrets является ошибкой. \\
\texttt{lean-build} & Показывает, что Lean принимает код с зарегистрированным toolchain. Не доказывает novelty/meaning. \\
\texttt{axiom-audit} & Делает transitive trusted-base traversal и применяет project axiom allowlist. \\
\texttt{exact-dependencies} & Строит elaborated dependency evidence и environment artifacts. \\
\bottomrule
\end{longtable}

PR mergeable только когда required machine checks зелёные \textbf{и} требуемые human reviews удовлетворены.

\subsection{Локальный pre-push gate}

\begin{lstlisting}
pytest -q
researchctl acceptance
researchctl agent-manifest
researchctl context-profiles
researchctl validate
researchctl workflow-audit
lake build
researchctl axioms-all
researchctl deps-exact-all
\end{lstlisting}

Для изменённого H-ID:

\begin{lstlisting}
researchctl semantic-diff H-0042
researchctl fingerprint H-0042
researchctl report H-0042
researchctl audit H-0042
\end{lstlisting}

\section{CD: scientific delivery и immutable release}

В обычном software CI/CD delivery означает deploy. Здесь delivery означает выпуск неизменяемого, воспроизводимого research artifact.

Перед release:

\begin{lstlisting}
researchctl validate
researchctl audit H-0042
researchctl fingerprint H-0042
researchctl packet H-0042
\end{lstlisting}

Verification packet должен содержать hypothesis card, Lean source, statement snapshot, proof reconstruction, literature/provenance evidence, reviews, dependency reports, environment lock, canonical agent prompt files, visibility profiles, blind-context input manifests, agent protocol manifest, provenance записей конкретных agent runs и SHA-256 manifest.

Research-release workflow запускается только для \status{VERIFIED\_RESULT} и должен быть защищён GitHub Environment с human approval. Рекомендуемый immutable tag:

\begin{lstlisting}
result/H-0042/<statement-hash-prefix>
\end{lstlisting}

Опубликованный tag не переписывается. Материальное изменение evidence создаёт новый release с явным описанием связи со старым.

\section{Merge policy для разных научных исходов}

\textbf{Successful theorem:} merge на соответствующей стадии, без прыжка сразу в \status{VERIFIED\_RESULT}.\\
\textbf{Refuted hypothesis:} merge воспроизводимого counterexample/refutation и статус \status{REFUTED}.\\
\textbf{Known result:} merge provenance/literature finding и статус \status{KNOWN\_RESULT}.\\
\textbf{Specification bug:} сохранить факт ошибочного encoding и поставить \status{SPECIFICATION\_BUG}; не переписывать историю.\\
\textbf{Blocked/inconclusive:} документировать причину и условие разблокировки.

\section{Правила работы агентов}

Перед proof-agent передавайте H-ID, frozen statement или явный \status{SPEC\_DRAFT}, список definitions которые можно/нельзя менять, axiom policy, имеющийся provenance/literature evidence и прямой запрет silent assumption strengthening.

Требуйте от proof-agent output в порядке: (1) interpretation exact statement; (2) proof architecture; (3) новые lemmas; (4) Lean implementation; (5) dependency/provenance notes; (6) явный запрос на specification change, если proof требует изменения постановки.

Literature-agent обязан давать reproducible queries и candidate prior art, а не бинарный novelty verdict. Skeptic-agent получает задачу искать counterexamples, удалять assumptions, проверять degenerate cases и искать stronger known results.

\begin{dangerbox}
Агенту запрещено: менять frozen statement без SCR; самостоятельно ставить \status{VERIFIED\_RESULT}; подписывать human reviews; придумывать citations/search logs; скрывать неудачу proof через усиление assumptions; трактовать green CI как научное одобрение.
\end{dangerbox}

\subsection{Каноническая структура prompt'ов и разделение ролей}

Нормативные prompt'ы агентов хранятся не в одном общем документе, а в каталоге \texttt{agents/}. Это часть воспроизводимой исследовательской среды.

\begin{lstlisting}
agents/
  common/
    SYSTEM.md
    OUTPUT_CONTRACT.md
  formalizer/
    SYSTEM.md
    TASK.md
    OUTPUT_SCHEMA.yaml
  skeptic/
  semantic-reviewer/
  proof-reviewer/
  literature-reviewer/
  provenance-auditor/
  reproducer/
  red-team/
  merge-gate/
\end{lstlisting}

Каждая роль получает общий system contract, свой role-specific system prompt, task template и machine-readable output schema. Один универсальный агент не должен одновременно генерировать гипотезу, доказывать её, проверять семантику, подтверждать novelty и одобрять merge: разделение ролей является частью защиты от self-review.

\begin{tabularx}{\textwidth}{@{}p{.27\textwidth}Y@{}}
\toprule
Роль & Основная обязанность \\
\midrule
\texttt{formalizer} & Формализовать ровно frozen statement; при необходимости semantic change остановиться и запросить SCR. \\
\texttt{skeptic} & Пытаться опровергнуть claim, находить boundary cases, лишние assumptions и trivialization. \\
\texttt{semantic-reviewer} & Сопоставить Lean definitions/statement с intended mathematics. \\
\texttt{proof-reviewer} & Восстановить человеческое доказательство и минимальный review surface. \\
\texttt{literature-reviewer} & Фальсифицировать novelty через воспроизводимый prior-art search. \\
\texttt{provenance-auditor} & Классифицировать происхождение содержательных зависимостей. \\
\texttt{reproducer} & Независимо воспроизвести build/evidence из чистой среды. \\
\texttt{red-team} & Автоматизированно искать блокирующие риски; не имеет права выдавать научное одобрение. \\
\texttt{merge-gate} & Проверить только repository evidence и policy gates перед merge. \\
\bottomrule
\end{tabularx}

\subsection{Хеширование prompt-протокола и provenance запуска}

Перед серьёзной AI-assisted работой создайте cryptographic manifest:

\begin{lstlisting}
researchctl agent-manifest
\end{lstlisting}

Команда вычисляет общий \texttt{agent\_protocol\_sha256} и отдельный hash каждой роли. Этот hash входит в research fingerprint и verification packet. Изменение канонического prompt'а поэтому является наблюдаемым изменением исследовательской среды.

Для подготовки prompt'а конкретной роли:

\begin{lstlisting}
researchctl agent-prompt skeptic H-0042
researchctl agent-prompt formalizer H-0042
\end{lstlisting}

Для фиксации реального AI-запуска:

\begin{lstlisting}
researchctl agent-run-init formalizer H-0042 \
  --model gpt-5.6-sol
\end{lstlisting}

Создаваемая запись обязана содержать H-ID, роль, модель, commit, frozen statement hash, общий protocol hash и role prompt hash. После завершения запуска человек или automation заполняет реальные input/output artifacts и outcome. Нельзя переписывать hash в старой записи для ``освежения'' evidence; при изменении prompt'а или statement создаётся новый run record.

\begin{warningbox}
\texttt{docs/PROMPTS.md} является только кратким индексом. При любом расхождении нормативными для агентов считаются файлы в \texttt{agents/}. Human sign-off при этом по-прежнему не может быть заменён agent output schema.
\end{warningbox}


\subsection{Blind Agent Context Layer: изоляция контекста}

Простого разделения ролей недостаточно. Если агенту дать весь checkout и только написать ``не смотри на другие этапы'', он всё равно знает глобальную цель, результаты других ролей и желаемый исход. Для независимых стадий используется \textbf{context firewall}: агент получает физически отдельный минимальный пакет, а не репозиторий.

\begin{principle}
Независимый agent run должен быть воспроизводим как функция от явного input packet. Отсутствие доступа к скрытому контексту обеспечивается структурой входа, а не обещанием агента игнорировать файлы.
\end{principle}

Канонические visibility profiles лежат в \texttt{agents/profiles/}.

\begin{lstlisting}
researchctl context-profiles
researchctl context-build skeptic H-0042 \
  --profile skeptic-blind
researchctl context-build proof-reviewer H-0042 \
  --profile proof-blind
\end{lstlisting}

Пакет создаётся в \texttt{reports/agent-inputs/H-XXXX/<RUN-ID>/} и содержит только разрешённые артефакты, нейтральный task и \texttt{INPUT\_MANIFEST.json}. Manifest фиксирует profile hash, statement hash, protocol hash, included artifacts, declared exclusions, sanitization policy, leakage scan и \texttt{context\_packet\_sha256}.

\begin{tabularx}{\textwidth}{@{}p{.25\textwidth}Y@{}}
\toprule
Уровень изоляции & Что скрывается \\
\midrule
Role separation & Только ответственность; общий контекст может быть известен. Слабая изоляция. \\
Output isolation & Outputs других агентов. \\
Task blindness & Желаемый глобальный исход, proof status, downstream цели. \\
Provenance blindness & Кто сгенерировал claim, generator reasoning, другие reviews. \\
Goal blindness & Максимально нейтральная локальная математическая постановка без project-specific цели. \\
\bottomrule
\end{tabularx}

Для \texttt{skeptic-blind} по умолчанию разрешены frozen statement, definitions, assumptions и notation; запрещены existing proof, proof reconstruction, Lean-success status, literature results, novelty claim, intended application и outputs других агентов. \texttt{proof-blind} не должен знать даже о существовании accepted formal proof: его задача --- независимо вывести человеческий proof sketch из statement/definitions/assumptions.

\subsubsection{Sanitization и пределы blindness}

Профиль может скрыть H-ID, заменить project-specific identifiers нейтральными alias'ами и удалить configured markers. Это снижает anchoring, но не даёт логической гарантии, что модель не догадается о тематике по самой математике. Корректная формулировка аудита: ``глобальная задача не была явно включена в зафиксированный input packet''.

\begin{dangerbox}
Blind-run не считается независимым, если агент имел полный checkout, доступ к git history/search, outputs других ролей или project overview. Для blind-run не используйте symlink на repository и не давайте shell с рабочей директорией проекта. Передавайте только скомпилированный context packet.
\end{dangerbox}

\subsubsection{Привязка запуска к пакету}

\begin{lstlisting}
researchctl agent-run-init skeptic H-0042 \
  --model <MODEL> \
  --profile skeptic-blind
\end{lstlisting}

Run record содержит visibility profile, путь к context packet, его SHA-256 и результат leakage scan. При изменении profile, statement или разрешённых входов создаётся новый packet/run; старые hashes не переписываются.

GitHub job \texttt{blind-context} строит smoke-packets для нескольких blind roles и выполняет автоматический leakage scan. Этот gate является required check вместе с \texttt{agent-protocol}. После merge новой версии обновите уже активный server-side ruleset командой \cmd{scripts/update\_github\_ruleset.sh OWNER/REPO}; bootstrap-скрипт существующий ruleset намеренно не перезаписывает.

\section{Incident response и sabotage testing}

Control system считается проверенным не одним зелёным build, а способностью правильно падать при атаке. После изменения infrastructure запускайте \cmd{researchctl acceptance} и полный matrix из \texttt{docs/FIRST\_RUN.md}.

Минимальные sabotage scenarios:

\begin{tabularx}{\textwidth}{@{}p{.36\textwidth}Y@{}}
\toprule
Саботаж & Ожидаемое поведение \\
\midrule
Добавить \texttt{sorry} & Lean может build'иться, но protocol gate обязан FAIL. \\
Изменить frozen binder/assumption & semantic diff показывает drift; validate FAIL. \\
Добавить custom axiom & fresh axiom audit находит её; policy FAIL без explicit allow. \\
Скрыть axiom через bridge lemma & transitive traversal всё равно находит axiom. \\
Изменить только proof & statement остаётся VALID; proof/reproduction reviews становятся STALE. \\
Сломать lean\_file/module/declaration & validate выдаёт controlled INVALID\_LEAN\_REFERENCE, без traceback. \\
Попытаться release до VERIFIED\_RESULT & CD workflow отказывает. \\
\bottomrule
\end{tabularx}

\section{Definition of Done для publishable result}

\begin{checklistbox}
\begin{itemize}
\item exact mathematical claim frozen;
\item no placeholder proofs;
\item clean Lean build;
\item trusted-base/axiom audit passes;
\item semantic review valid;
\item human proof reconstruction complete;
\item proof review valid;
\item key dependency provenance classified;
\item meaningful assumptions challenged/ablated where appropriate;
\item literature search reproducibly logged;
\item novelty review valid and claim language calibrated;
\item independent reproduction valid;
\item exact dependency report current;
\item environment lock current;
\item final review valid;
\item required distinct human signers recorded;
\item verification packet generated;
\item immutable release created from reviewed commit.
\end{itemize}
\end{checklistbox}

Это и есть intended meaning \status{VERIFIED\_RESULT}.

\appendix
\section{Краткая последовательность для новой гипотезы}

\begin{enumerate}
\item \cmd{researchctl new H-XXXX "..."}; создать Issue и Draft PR.
\item Заполнить intended statement, scope, definitions, assumptions и initial prior art.
\item На \status{SPEC\_DRAFT} провести semantic/falsification work.
\item \cmd{researchctl validate}; \cmd{researchctl freeze H-XXXX}; проверить \cmd{semantic-diff}.
\item Разрабатывать proof, не меняя frozen statement скрытно.
\item Пройти build, placeholder, axiom и exact dependency checks.
\item Выполнить semantic review.
\item Восстановить proof человеком и выполнить proof review.
\item Выполнить ablation/adversarial tests.
\item Выполнить literature/provenance review с reproducible search log.
\item Независимо воспроизвести на clean checkout.
\item Проверить freshness всех reviews и final claim language.
\item Создать verification packet и immutable research release.
\end{enumerate}

\section{Карта файлов}

\begin{tabularx}{\textwidth}{@{}p{.38\textwidth}Y@{}}
\toprule
Путь & Назначение \\
\midrule
\texttt{hypotheses/H-XXXX.yaml} & machine-readable scientific state \\
\texttt{MathCI/*.lean} & formal source \\
\texttt{proofs/H-XXXX.md} & рабочая human proof reconstruction \\
\texttt{lit/H-XXXX.md} & reproducible literature search log \\
\texttt{provenance/H-XXXX.yaml} & dependency provenance registry \\
\texttt{reviews/H-XXXX/*.yaml} & hash-bound review records \\
\texttt{reports/statements/} & frozen statement snapshots \\
\texttt{reports/dependencies/} & source/exact dependency evidence \\
\texttt{reports/fingerprints/} & domain/evidence fingerprints \\
\texttt{research-policy.yaml} & policy-as-code \\
\texttt{docs/HUMAN\_MANUAL.tex} & canonical human SOP \\
\texttt{AGENTS.md} & concise machine/agent contract \\
\texttt{agents/} & canonical versioned prompts, visibility profiles and output schemas for AI roles \\
\texttt{reports/agent-inputs/H-XXXX/} & isolated blind-context packets with input manifests \\
\texttt{agent-runs/H-XXXX/} & provenance records of concrete AI-assisted runs \\
\texttt{reports/agents/manifest.json} & cryptographic manifest of the agent protocol \\
\bottomrule
\end{tabularx}

\section{Команды ежедневного контроля}

\begin{lstlisting}
# Repository baseline
pytest -q
researchctl acceptance
researchctl agent-manifest
researchctl validate
researchctl workflow-audit
lake build
researchctl axioms-all
researchctl deps-exact-all

# One hypothesis
researchctl semantic-diff H-0042
researchctl fingerprint H-0042
researchctl report H-0042
researchctl audit H-0042

# Evidence lifecycle
researchctl provenance-init H-0042
researchctl ablate H-0042
researchctl review-init H-0042
researchctl packet H-0042
\end{lstlisting}

\vfill
\begin{center}
\small\textcolor{gray}{Этот документ задаёт human operating procedure. При расхождении краткого Markdown summary и данного TeX manual приоритет имеет TeX manual, если policy-as-code/CI не устанавливают более строгого требования.}
\end{center}
\end{document}
